\documentclass[11pt]{article}

\usepackage[margin=1in]{geometry}
\usepackage{amsmath, amssymb}

\title{\textbf{Week 3 Writing Assignment (CS-EE5841)}\\Convolutional Neural Networks and Advanced CNN Architectures}
\author{}
\date{}

\begin{document}

\maketitle

\noindent \textbf{Total Points: 50}

\vspace{0.3cm}

\noindent \textbf{Instructions:} Show all key steps using formulas. Provide brief interpretation where needed.

%--------------------------------------------------
\section*{Q1. Convolution Output Size and Parameter Count (12 pts)}

\subsection*{Question}

An input image has size:
\[
32 \times 32 \times 3
\]

A convolutional layer uses:
\[
16 \text{ filters}, \quad 5 \times 5 \text{ kernel}, \quad S=1, \quad P=2
\]

Compute:
\begin{itemize}
    \item Output spatial dimension
    \item Output volume size
    \item Number of trainable parameters
\end{itemize}

\subsection*{Hint}

\[
O=\frac{I-F+2P}{S}+1
\]

\[
\text{Parameters}=(F_hF_wC_{in}+1)C_{out}
\]

\subsection*{Solution}

--- Write solutions for all tasks here 

\noindent \textbf{Interpretation:} ... also write brief interpretation.

%--------------------------------------------------
\section*{Q2. Receptive Field Calculation in a CNN (10 pts)}

\subsection*{Question}

Consider a CNN with three convolutional layers. Each layer uses:
\[
3 \times 3 \text{ kernels}, \quad S=1, \quad P=1
\]

Compute the receptive field size after each convolutional layer.

\subsection*{Hint}

For stride 1 convolution:

\[
R_l=R_{l-1}+(K_l-1)
\]

where:
\[
R_0=1
\]

\subsection*{Solution}

--- Write solutions for all tasks here 

\noindent \textbf{Interpretation:} ... also write brief interpretation.

%--------------------------------------------------
\section*{Q3. Max Pooling and Feature Map Reduction (8 pts)}

\subsection*{Question}

Given the feature map:

\[
X=
\begin{bmatrix}
1 & 3 & 2 & 4\\
5 & 6 & 1 & 2\\
7 & 2 & 8 & 3\\
4 & 1 & 5 & 9
\end{bmatrix}
\]

Apply \(2 \times 2\) max pooling with stride \(S=2\).

\subsection*{Solution}

--- Write solutions for all tasks here 

\noindent \textbf{Interpretation:} ... also write brief interpretation.

%--------------------------------------------------
\section*{Q4. Residual Learning in ResNet (10 pts)}

\subsection*{Question}

A residual block computes:

\[
y=F(x)+x
\]

Suppose:

\[
x=
\begin{bmatrix}
1\\
-2\\
3
\end{bmatrix},
\quad
F(x)=
\begin{bmatrix}
0.5\\
1.0\\
-1.5
\end{bmatrix}
\]

Compute the residual block output \(y\). Then briefly explain why this formulation helps train very deep CNNs.

\subsection*{Solution}

--- Write solutions for all tasks here 

\noindent \textbf{Interpretation:} ... also write brief interpretation.

%--------------------------------------------------
\section*{Q5. U-Net Skip Connections and DenseNet Feature Reuse (10 pts)}

\subsection*{Question}

Answer the following conceptual-mathematical questions.

\begin{enumerate}
    \item In U-Net, suppose the decoder feature map has 128 channels and the corresponding encoder skip feature map has 128 channels. If they are concatenated, how many channels enter the next convolutional block?
    \item In DenseNet, suppose a dense block starts with 64 channels and has growth rate \(k=32\). After 4 dense layers, how many channels are available?
\end{enumerate}

\subsection*{Hint}

For concatenation:

\[
C_{out}=C_1+C_2
\]

For DenseNet:

\[
C_L=C_0+kL
\]

\subsection*{Solution}

--- Write solutions for all tasks here 

\noindent \textbf{Interpretation:} ... also write brief interpretation.

\end{document}